% !PRINTED 370
\weberpage{370}
% Source-faithful three-stroke Cayley configuration sign.
\providecommand{\weberthreeup}{\mathord{\tikz[baseline=-0.2ex,x=0.23em,y=0.23em,line width=0.45pt]{\draw (0,0)--(-1,1) (0,0)--(0,1.08) (0,0)--(1,1);}}}
%<<CONT
enthält, und der mit (1) ein azygetisches Paar bildet. Er muss
also von der Gestalt sein:
\weq{2}{x_1 y_2,\ x_3 z_3,\ x_4 z_4,\ x_5 z_5,\ x_6 z_6,\ x_7 z_7,}
und demnach ist $y_2$ auch durch $[2\,1]$ zu bezeichnen, w. z. b. w.

Die $z$ in (2) sind von den $y$ in (1) verschieden, und da z. B.
$z_3 = [2\,3]$ ist, so folgt, dass allgemein zwei $[\mu\,\nu]$, die nicht in
beiden Ziffern $\mu, \nu$ übereinstimmen, von einander verschieden
sind. Da man aus sieben Ziffern einundzwanzig Paare bilden
kann, so erhält man auf diese Weise alle Doppeltangenten und
jede nur einmal.

Aus dieser Darstellungsweise ergiebt sich auch, dass keine
azygetischen Systeme von mehr als sieben Doppeltangenten existiren
; denn fügen wir zu den sieben $x$ noch eine beliebige
weitere Doppeltangente, für die wir bei der Gleichberechtigung
der Ziffern 1 bis 7 etwa $y_2 = [1\,2]$ wählen können, so ist $y_2 x_1 x_2$
ein syzygetisches Tripel.

Es ist zweckmässig, eine achte Ziffer 8 einzuführen, und die
Elemente des ursprünglichen Siebener-Systems nicht durch die
einfachen Ziffern, sondern durch die Paare
\[
[1\,8],\ [2\,8],\ [3\,8],\ [4\,8],\ [5\,8],\ [6\,8],\ [7\,8]
\]
zu bezeichnen, wobei dann auch gelten soll, dass $[1\,8] = [8\,1]$ u. s. f.
ist. Dann werden alle 28 Doppeltangenten übereinstimmend
durch die Paare $[\mu\,\nu]$ bezeichnet, in denen $\mu$ und $\nu$ zwei verschiedene
 Ziffern der Reihe 1 bis 8 bedeuten.

Dabei ist es für die Uebersichtlichkeit sehr förderlich, wenn
man eine anschauliche Bezeichnung anwendet\footnote{Nach \gesperrt{Cayley}; vergl. \gesperrt{Salmon}, \glqq Higher plane curves\grqq{}.}. Man deutet
eine Doppeltangente $[\mu\,\nu]$ durch einen einfachen Strich $\mid$ an, an
dessen Enden man sich die beiden Ziffern $\mu, \nu$ gesetzt denkt.
Dann bedeuten zwei Striche ohne gemeinsamen Punkt $\mid\mid$ zwei
Doppeltangenten, in deren Bezeichnung $[\mu\,\nu]$ keine gemeinschaftliche
 Ziffer vorkommt, und zwei von einem Punkte auslaufende
Striche, $\vee$, zwei Doppeltangenten, die in ihren Symbolen $[\mu\,\nu]$
eine gemeinschaftliche Ziffer haben. Hiernach sind die complicirteren
 Zeichen, die wir nachher anwenden, von selbst verständlich.
Für ein Tripel von Doppeltangenten haben wir z. B. folgende
fünf Zeichen $\mid\mid\mid$, $\sqcap$, $\triangle$, $\weberthreeup$, $\vee\mid$.
